\chapter{Transurethral Resection of Bladder Tumor}

\section*{Introduction \& Clinical Indications}
Transurethral resection of bladder tumor (TURBT) removes visible bladder lesions through a cystoscope passed along the urethra. It provides tissue for diagnosis, grading, and staging and is a key first procedure for suspected non-muscle-invasive bladder cancer. A complete, well-oriented resection helps determine whether detrusor muscle is present and whether repeat resection or additional treatment is needed.

TURBT is not the same as radical cystectomy. The treatment plan depends on tumor number, size, appearance, grade, depth, carcinoma in situ, recurrence pattern, and the patient’s overall risk. Selected patients receive immediate intravesical chemotherapy or later bladder-directed treatment; high-risk or muscle-invasive disease may require a different pathway.

\section*{Relevant Surgical Anatomy}
The urethra provides the endoscopic route to the bladder. The trigone, ureteric orifices, bladder neck, lateral walls, dome, and posterior wall must be mapped because location affects resection safety and risk to the ureters or obturator nerve. Thin bladder walls at the dome and the obturator reflex near the lateral wall are important technical considerations.

\section*{Preoperative Workup \& Preparation}
Evaluate hematuria, urine culture, renal function, blood count, anticoagulants, imaging, and prior pathology. Treat active urinary infection when possible and plan management of antithrombotic therapy. Explain bleeding, infection, dysuria, urinary retention, ureteric injury, bladder perforation, incomplete sampling, need for repeat TURBT, intravesical therapy, or radical surgery. Consent should address the possibility of a catheter after the procedure.

\section*{Surgical Technique \& Procedure}
Under suitable anesthesia, cystoscopy examines the entire bladder and documents lesion number, size, location, and appearance. Tumor is removed in a way that preserves orientation and obtains the appropriate deep tissue, with separate specimens when useful. Hemostasis is achieved, the bladder is drained, and a catheter is placed if indicated. A single postoperative intravesical chemotherapy dose may be considered for selected low- or intermediate-risk cases when there is no suspected perforation or significant bleeding.

\section*{Complications \& Risks}
Risks include hematuria, clot retention, urinary infection, urethral injury, bladder perforation, obturator jerk, ureteric injury, incomplete resection, and anesthesia complications. Delayed complications include stricture, recurrent tumors, progression, and repeated-procedure burden. Severe bleeding, fever, worsening abdominal pain, inability to urinate, or reduced urine output requires prompt review.

\section*{Postoperative Care \& Recovery}
Monitor urine color, catheter drainage, pain, fever, and bladder spasms. Catheter irrigation may be needed for clots. Review pathology with the patient and document the next step, which may include repeat TURBT, intravesical therapy, surveillance cystoscopy, systemic therapy, or referral for radical treatment. Smoking cessation and reliable surveillance are important because recurrence is common.

\section*{Outcomes \& Prognosis}
TURBT is diagnostic and therapeutic for many non-muscle-invasive tumors, but recurrence and progression risk vary widely. Prognosis depends on depth, grade, carcinoma in situ, tumor burden, response to intravesical treatment, and adherence to surveillance rather than on the endoscopic procedure alone.

\section*{References}
\begin{enumerate}
\item European Association of Urology. Guidelines on Non-Muscle-Invasive Bladder Cancer.
\item American Urological Association/Society of Urologic Oncology. Non-Muscle-Invasive Bladder Cancer Guideline.
\item National Institute for Health and Care Excellence. Bladder cancer: diagnosis and management.
\end{enumerate}
